% Queueing diagram: how processes move between the ready queue, the CPU
% and I/O queues.
% Usage: % Queueing diagram: how processes move between the ready queue, the CPU
% and I/O queues.
% Usage: % Queueing diagram: how processes move between the ready queue, the CPU
% and I/O queues.
% Usage: % Queueing diagram: how processes move between the ready queue, the CPU
% and I/O queues.
% Usage: \input{figures/l02-queues}   (width ~108 mm)
\begin{tikzpicture}[x=1mm, y=1mm, font=\footnotesize\sffamily,
  >={Stealth[length=1.6mm]},
  res/.style={draw, circle, fill=ln-accent!22, minimum size=10mm, inner sep=0pt},
  bx/.style={draw, rounded corners=2pt, fill=white, minimum height=6mm,
             minimum width=22mm, inner sep=1pt, font=\scriptsize\sffamily},
  lab/.style={font=\scriptsize\sffamily}]
  % queue drawing: rectangle with slots, open at the left
  \newcommand{\lnqueue}[3]{% x, y, label
    \draw[fill=ln-box] (#1,#2-3) rectangle ++(28,6);
    \foreach \i in {1,...,4} \draw (#1+28-4*\i,#2-3) -- ++(0,6);
    \node[lab, anchor=south] at (#1+14,#2+3) {#3};}
  \lnqueue{16}{40}{\iflangja{レディキュー}{ready queue}}
  \node[res] (cpu) at (64,40) {CPU};
  \draw[->] (44,40) -- (cpu.west);
  % right-hand bus and left-hand return line
  \draw (cpu.east) -- (102,40) -- (102,-8);
  \draw[->] (2,-8) -- (2,40) -- (16,40);
  % I/O
  \draw[->] (102,26) -- node[lab, above] {\iflangja{I/O 要求}{I/O request}} (80,26);
  \lnqueue{52}{26}{\iflangja{I/O キュー}{I/O queue}}
  \node[res, minimum size=8mm] (io) at (36,26) {I/O};
  \draw[->] (52,26) -- (io.east);
  \draw (io.west) -- (2,26);
  % timeslice expired
  \draw (102,14) -- node[lab, above, pos=0.5] {\iflangja{タイムスライス満了}{timeslice expired}} (2,14);
  % fork
  \draw[->] (102,3) -- node[lab, above] {\iflangja{子を fork}{fork a child}} (74,3);
  \node[bx] (ch) at (63,3) {\iflangja{子プロセスが作られる}{child is created}};
  \draw (ch.west) -- (2,3);
  % interrupt
  \draw[->] (102,-8) -- node[lab, above] {\iflangja{割り込みを待つ}{wait for an interrupt}} (74,-8);
  \node[bx] (in) at (63,-8) {\iflangja{割り込みが発生}{interrupt occurs}};
  \draw (in.west) -- (2,-8);
\end{tikzpicture}
   (width ~108 mm)
\begin{tikzpicture}[x=1mm, y=1mm, font=\footnotesize\sffamily,
  >={Stealth[length=1.6mm]},
  res/.style={draw, circle, fill=ln-accent!22, minimum size=10mm, inner sep=0pt},
  bx/.style={draw, rounded corners=2pt, fill=white, minimum height=6mm,
             minimum width=22mm, inner sep=1pt, font=\scriptsize\sffamily},
  lab/.style={font=\scriptsize\sffamily}]
  % queue drawing: rectangle with slots, open at the left
  \newcommand{\lnqueue}[3]{% x, y, label
    \draw[fill=ln-box] (#1,#2-3) rectangle ++(28,6);
    \foreach \i in {1,...,4} \draw (#1+28-4*\i,#2-3) -- ++(0,6);
    \node[lab, anchor=south] at (#1+14,#2+3) {#3};}
  \lnqueue{16}{40}{\iflangja{レディキュー}{ready queue}}
  \node[res] (cpu) at (64,40) {CPU};
  \draw[->] (44,40) -- (cpu.west);
  % right-hand bus and left-hand return line
  \draw (cpu.east) -- (102,40) -- (102,-8);
  \draw[->] (2,-8) -- (2,40) -- (16,40);
  % I/O
  \draw[->] (102,26) -- node[lab, above] {\iflangja{I/O 要求}{I/O request}} (80,26);
  \lnqueue{52}{26}{\iflangja{I/O キュー}{I/O queue}}
  \node[res, minimum size=8mm] (io) at (36,26) {I/O};
  \draw[->] (52,26) -- (io.east);
  \draw (io.west) -- (2,26);
  % timeslice expired
  \draw (102,14) -- node[lab, above, pos=0.5] {\iflangja{タイムスライス満了}{timeslice expired}} (2,14);
  % fork
  \draw[->] (102,3) -- node[lab, above] {\iflangja{子を fork}{fork a child}} (74,3);
  \node[bx] (ch) at (63,3) {\iflangja{子プロセスが作られる}{child is created}};
  \draw (ch.west) -- (2,3);
  % interrupt
  \draw[->] (102,-8) -- node[lab, above] {\iflangja{割り込みを待つ}{wait for an interrupt}} (74,-8);
  \node[bx] (in) at (63,-8) {\iflangja{割り込みが発生}{interrupt occurs}};
  \draw (in.west) -- (2,-8);
\end{tikzpicture}
   (width ~108 mm)
\begin{tikzpicture}[x=1mm, y=1mm, font=\footnotesize\sffamily,
  >={Stealth[length=1.6mm]},
  res/.style={draw, circle, fill=ln-accent!22, minimum size=10mm, inner sep=0pt},
  bx/.style={draw, rounded corners=2pt, fill=white, minimum height=6mm,
             minimum width=22mm, inner sep=1pt, font=\scriptsize\sffamily},
  lab/.style={font=\scriptsize\sffamily}]
  % queue drawing: rectangle with slots, open at the left
  \newcommand{\lnqueue}[3]{% x, y, label
    \draw[fill=ln-box] (#1,#2-3) rectangle ++(28,6);
    \foreach \i in {1,...,4} \draw (#1+28-4*\i,#2-3) -- ++(0,6);
    \node[lab, anchor=south] at (#1+14,#2+3) {#3};}
  \lnqueue{16}{40}{\iflangja{レディキュー}{ready queue}}
  \node[res] (cpu) at (64,40) {CPU};
  \draw[->] (44,40) -- (cpu.west);
  % right-hand bus and left-hand return line
  \draw (cpu.east) -- (102,40) -- (102,-8);
  \draw[->] (2,-8) -- (2,40) -- (16,40);
  % I/O
  \draw[->] (102,26) -- node[lab, above] {\iflangja{I/O 要求}{I/O request}} (80,26);
  \lnqueue{52}{26}{\iflangja{I/O キュー}{I/O queue}}
  \node[res, minimum size=8mm] (io) at (36,26) {I/O};
  \draw[->] (52,26) -- (io.east);
  \draw (io.west) -- (2,26);
  % timeslice expired
  \draw (102,14) -- node[lab, above, pos=0.5] {\iflangja{タイムスライス満了}{timeslice expired}} (2,14);
  % fork
  \draw[->] (102,3) -- node[lab, above] {\iflangja{子を fork}{fork a child}} (74,3);
  \node[bx] (ch) at (63,3) {\iflangja{子プロセスが作られる}{child is created}};
  \draw (ch.west) -- (2,3);
  % interrupt
  \draw[->] (102,-8) -- node[lab, above] {\iflangja{割り込みを待つ}{wait for an interrupt}} (74,-8);
  \node[bx] (in) at (63,-8) {\iflangja{割り込みが発生}{interrupt occurs}};
  \draw (in.west) -- (2,-8);
\end{tikzpicture}
   (width ~108 mm)
\begin{tikzpicture}[x=1mm, y=1mm, font=\footnotesize\sffamily,
  >={Stealth[length=1.6mm]},
  res/.style={draw, circle, fill=ln-accent!22, minimum size=10mm, inner sep=0pt},
  bx/.style={draw, rounded corners=2pt, fill=white, minimum height=6mm,
             minimum width=22mm, inner sep=1pt, font=\scriptsize\sffamily},
  lab/.style={font=\scriptsize\sffamily}]
  % queue drawing: rectangle with slots, open at the left
  \newcommand{\lnqueue}[3]{% x, y, label
    \draw[fill=ln-box] (#1,#2-3) rectangle ++(28,6);
    \foreach \i in {1,...,4} \draw (#1+28-4*\i,#2-3) -- ++(0,6);
    \node[lab, anchor=south] at (#1+14,#2+3) {#3};}
  \lnqueue{16}{40}{\iflangja{レディキュー}{ready queue}}
  \node[res] (cpu) at (64,40) {CPU};
  \draw[->] (44,40) -- (cpu.west);
  % right-hand bus and left-hand return line
  \draw (cpu.east) -- (102,40) -- (102,-8);
  \draw[->] (2,-8) -- (2,40) -- (16,40);
  % I/O
  \draw[->] (102,26) -- node[lab, above] {\iflangja{I/O 要求}{I/O request}} (80,26);
  \lnqueue{52}{26}{\iflangja{I/O キュー}{I/O queue}}
  \node[res, minimum size=8mm] (io) at (36,26) {I/O};
  \draw[->] (52,26) -- (io.east);
  \draw (io.west) -- (2,26);
  % timeslice expired
  \draw (102,14) -- node[lab, above, pos=0.5] {\iflangja{タイムスライス満了}{timeslice expired}} (2,14);
  % fork
  \draw[->] (102,3) -- node[lab, above] {\iflangja{子を fork}{fork a child}} (74,3);
  \node[bx] (ch) at (63,3) {\iflangja{子プロセスが作られる}{child is created}};
  \draw (ch.west) -- (2,3);
  % interrupt
  \draw[->] (102,-8) -- node[lab, above] {\iflangja{割り込みを待つ}{wait for an interrupt}} (74,-8);
  \node[bx] (in) at (63,-8) {\iflangja{割り込みが発生}{interrupt occurs}};
  \draw (in.west) -- (2,-8);
\end{tikzpicture}
